\section{Heat-potential contacts and separated configurations}
\label{sec:contact}

The contact construction replaces a positive signed integral by a positive
measure carried by selected space--scale points.  Its useful property is
an initial sign and a cumulative sign on each outward vertical.  We first
give this construction without any geometric assumptions.  We then use
it, together with the elementary estimates of \Cref{sec:local}, to
reduce a separated configuration to smaller multiplicities.

The extension viewpoint of Caffarelli and Silvestre
\cite[Sections~1--2]{CS2007} motivated an earlier potential-based approach
to the contact construction.  The present construction uses a heat
potential and is proved locally; their fractional-Laplacian extension
theorem is not invoked as a varifold regularity input.

\subsection{The entrance boundary and its contact measure}

Let $X$ be a flat torus, with its volume measure denoted by $\dd y$.
Gaussian kernels on $X$ are periodized, and preserve constants.  Fix a
bottom scale $a>0$.  A \emph{continuation set} is a closed subset
$\mathcal D\subset X\times[a,\infty)$ such that
\[
 (x,s)\in\mathcal D\ \Longrightarrow\ (x,v)\in\mathcal D
 \quad(v\ge s),\qquad
 X\times[s_{\rm top},\infty)\subset\mathcal D
\]
for some $s_{\rm top}>a$.  Its entrance boundary is
$K=\mathcal D\setminus\operatorname{int}\mathcal D$, where the interior is
taken in $X\times(0,\infty)$.  Thus $K$ is compact.  No artificial upper
time boundary belongs to $K$.

Here are the two types of continuation sets used below.  For a finite
mesh $t_j\in[a,s_{\rm top})$ and continuous tests $q_j$, the inequalities
$q_j(x)\le\lambda_j$ are required only at strictly larger mesh scales:
\begin{equation}\label{ct:mesh-domain}
 \mathcal D=\{s\ge a\}\cap
 \bigcap_j\bigl(\{s\ge t_j\}\cup\{q_j(x)\le\lambda_j\}\bigr).
\end{equation}
Each factor is closed and upward closed.  At $s=t_j$ the current test is
omitted.  With tests continuous in the scale, impose the non-strict
continuation inequality at every scale $v\ge s$ at which a test is
defined, and extend by the full cylinder above the last test.  In our
applications all tests hold strictly near that last scale, so this
extension is closed.  Equivalently one can first impose the inequalities
at $v>s$ and pass to $v=s$ by continuity.  Taking a limit of points of
$\mathcal D$ proves closedness: any fixed $v>s$ is eventually larger
than the approximating scales.

The entrance boundary includes bottom points, current hits, and vertical
sides below an equality at a larger scale.  These sides are essential:
a tangential equality can prevent a spatial neighborhood from belonging
to the continuation set even though the outward vertical itself remains
in that set.  Upward closure implies that a backward space--time cylinder
contained in $\mathcal D$ provides forward room as well; hence the above
definition agrees with the parabolic entrance boundary.

\begin{lemma}[Contact measure]\label{ct:contact}
Let $f$ be a finite signed Borel measure on $X$, with $M_f=f(X)>0$, and
let $K$ be the entrance boundary of a continuation set.  There is a
probability measure $\sigma$ on $K$ such that the positive spatial measure
\[
 h=M_f(\pi_X)_\#\sigma
\]
has mass $M_f$ and, at every $(x,s_0)\in\spt\sigma$,
\begin{equation}\label{ct:contact-signs}
 (f-h)_{s_0}(x)\ge0,
 \qquad
 \int_{s_0}^{s}v^2(f-h)_v(x)\,\dd\log v\ge0
 \quad(s\ge s_0).
\end{equation}
The second inequality concerns the integral; it does not assert a
pointwise sign at scales larger than $s_0$.
\end{lemma}

\begin{proof}
For $\sigma\in\mathcal P(K)$ put
$g_\sigma=f-M_f(\pi_X)_\#\sigma$.  This measure has zero mass.  Define
\[
 b_\sigma(x,s)=G_s*(-\Delta_X)^{-1}g_\sigma(x),
\]
using the inverse with zero spatial mean.  If $\lambda_k$ are the
nonzero eigenvalues of $-\Delta_X$, its kernel has Fourier coefficients
$e^{-\lambda_k s^2/2}/\lambda_k$.  Since $s\ge a>0$, the kernel and all
its derivatives converge uniformly on bounded scale intervals.  Thus
$b_\sigma$ is smooth, and
\begin{equation}\label{ct:potential-derivative}
 \frac{\dd b_\sigma}{\dd\log s}=-s^2(g_\sigma)_s,
 \qquad \partial_t b_\sigma=\Delta_X b_\sigma,
 \qquad t=s^2/2.
\end{equation}
Weak convergence of probabilities implies uniform convergence of their
potentials on $K$: the kernel is jointly continuous on the relevant
compact product, and can be approximated uniformly by finite sums of
products of continuous functions of its two variables.

Choose finite $1/j$-nets $K_j=\{z_{j1},\ldots,z_{jN_j}\}$ in $K$.
For a probability vector $p$ on $K_j$ and $\beta>0$, the map
\[
 p\longmapsto
 \left(\frac{\exp(\beta b_p(z_{ji}))}
 {\sum_{k=1}^{N_j}\exp(\beta b_p(z_{jk}))}\right)_{i=1}^{N_j}
\]
is continuous from the probability simplex to itself.  Brouwer's Fixed
Point Theorem gives a fixed point \cite[Theorem~2]{Milnor1978}.  First let $\beta\to\infty$ on a
fixed net and take a convergent subsequence.  If a point's limiting
payoff is at least $\delta$ below the maximum, its fixed-point weight is
at most $e^{-\beta\delta/2}$ for all sufficiently large $\beta$.
The limiting probability is therefore carried by maximizers of its own
potential on $K_j$.  Next pass to a weak limit as $j\to\infty$.
Uniform convergence of the potentials and density of the nets show that
the resulting probability $\sigma$ satisfies
\begin{equation}\label{ct:self-maximizers}
 \spt\sigma\subset\operatorname*{arg\,max}_{z\in K}b_\sigma(z).
\end{equation}
Indeed a compact subset separated from the limiting maximum by $\delta$
receives no mass from sufficiently late net probabilities; a nearby net
point witnesses the strictly larger payoff.  Exhaust the complement of
the closed maximizer set by such compact subsets.

Write $b=b_\sigma$.  Its maximum on the whole continuation set is
bounded by $\max_K b$.  To verify this on an irregular boundary, truncate
at any scale $S>s_{\rm top}$ and apply the elementary maximum argument to
$b-\epsilon t$.  At a maximum outside $K$ below the artificial top,
$\partial_t(b-\epsilon t)-\Delta(b-\epsilon t)=-\epsilon$ contradicts
the derivative conditions.  At a maximum at the artificial top, the
one-sided time derivative is nonnegative and the spatial Laplacian is
nonpositive, giving the same contradiction.  Let $\epsilon\downarrow0$.
This argument uses only that points outside $K$ have an open
neighborhood in the continuation set, and hence also covers arbitrary
vertical-side geometry.

For $(x,s_0)\in\spt\sigma$, \eqref{ct:self-maximizers} and upward
closure now give $b(x,s_0)\ge b(x,s)$ for every $s\ge s_0$.
The right derivative at $s_0$ is nonpositive.  Applying
\eqref{ct:potential-derivative} proves the first inequality in
\eqref{ct:contact-signs}; integrating that identity proves the second.
\end{proof}

\begin{figure}[tb]
\centering
\begin{tikzpicture}[x=1.1cm,y=1.1cm,>=Latex,font=\small]
 \fill[blue!6] (-3,.75)--(0,.75)--(0,2.15)--(3,2.15)--(3,3.55)--(-3,3.55)--cycle;
 \draw[->] (-3.45,0)--(3.55,0) node[right] {$x$};
 \draw[->] (-3.45,0)--(-3.45,3.85) node[above] {$\log s$};
 \draw[blue!65!black,very thick] (-3,.75)--(0,.75)--(0,2.15)--(3,2.15);
 \draw[densely dashed,gray] (-3.35,3.3)--(3.1,3.3);
 \node[left] at (-3.45,3.3) {$\log\rho$};
 \node[blue!45!black] at (1.5,2.7) {continuation region};
 \node[anchor=east,align=right] at (-.5,1.9) {vertical side of the\\entrance boundary};
 \draw[->,gray] (-.48,1.7)--(-.08,1.6);
 \draw[orange!75!black,thick,->] (.05,1.3)--(.05,3.13);
 \fill (0,1.27) circle (2pt);
 \node[right] at (.22,1.22) {contact at $\underline s$};
 \draw[fill=white,blue!65!black,thick] (0,2.15) circle (2pt);
 \node[right,fill=white,inner sep=2pt] at (.25,2.02) {threshold equality at $s_0$};
\end{tikzpicture}
\caption{A schematic entrance boundary. In this illustration the contact
scale is denoted by $\underline s$ and a larger threshold scale by $s_0$.
A contact on a vertical side can have
derivative scale $\underline s<s_0$, where $s_0$ is a larger threshold scale.
The contact rule controls the initial derivative and the cumulative integral
along the outward vertical. It does not impose a pointwise sign at every later
scale. The two scales remain distinct in the final mass charge.}
\label{fig:contact}
\end{figure}


\subsection{Localization of the signed source}

All extensions in the next statement are made consistently: if $L'$ is
restricted to a working chart, $\tau$ is restricted to the same or a
larger chart, so that $|L'|\le C_0\tau$ survives extension by zero.
The constant density $d\,\dd y$ is placed on the entire torus.

\begin{lemma}[Uniform weight localization]\label{ct:localization}
Let $|L'|\le C_0\tau$, where $\tau$ is a finite positive measure on a
torus of diameter comparable to $R$, and let
$0\le w\le1$, $\Lip w\le R^{-1}$.  Put
\[
 \nu=\tau+d\,\dd y,\qquad f=wL'-\eta\nu,
 \qquad 0<\eta<1.
\]
Suppose, for some fixed $C_1$, that
$\tau(X)/(dR^m)\le e^{C_1T}$.  For all sufficiently large $T$, if
\begin{equation}\label{ct:localization-scale}
 0<\rho\le c\eta^2R/T^2,
\end{equation}
then, uniformly for every $x\in X$ and $0<s\le\rho$,
\begin{equation}\label{ct:localized-source}
 f_s(x)\le w(x)L'_s(x)-\frac\eta2\bigl(\tau_s(x)+d\bigr).
\end{equation}
Here $c>0$ depends only on $C_0,C_1$ and the fixed torus shape.
If $\sigma,h$ are furnished by \Cref{ct:contact} with contacts at
scales $s_0\le\rho$, then $w(x)>0$ at each contact, and
\begin{align}
 L'_{s_0}(x)&\ge h_{s_0}(x)+\tfrac\eta2\nu_{s_0}(x),
 \label{ct:localized-initial}\\
 \int_{s_0}^s v^2\bigl(L'-\tfrac\eta2\nu\bigr)_v(x)
 \,\dd\log v&\ge0\qquad(s_0\le s\le\rho).
 \label{ct:localized-cumulative}
\end{align}
In particular, contacts lie in the interior region containing
$\{w>0\}$.  These conclusions are uniform as the bottom scale tends to
zero.
\end{lemma}

\begin{proof}
The only commutator is
\[
 f_s(x)=w(x)L'_s(x)-\eta\nu_s(x)+E_s(x),\qquad
 |E_s(x)|\le C_0\int G_s(x-y)|w(y)-w(x)|\,\dd\tau(y).
\]
Split at torus distance $r_*=c_0\eta R$.  For sufficiently small fixed
$c_0$, the near contribution is at most $\eta\tau_s(x)/4$.
The far contribution is at most
\[
 C\tau(X)s^{-m}\exp(-c_2\eta^2R^2/s^2).
\]
Writing $x_*=s/R$, its ratio to $d$ is bounded by
$Ce^{C_1T}x_*^{-m}e^{-c_2\eta^2/x_*^2}$.
On $0<x_*\le c\eta^2/T^2$ this is $o(\eta)$ uniformly:
the negative exponent is at least a fixed multiple of $T^4/\eta^2$,
which dominates $C_1T$, $m|\log x_*|$, and $|\log\eta|$.
The same estimate applies to the nonnearest periodized copies.
This proves \eqref{ct:localized-source}.

At a contact $f_{s_0}\ge h_{s_0}\ge0$.  If $w(x)=0$,
\eqref{ct:localized-source} would give a strictly negative value.
Division by $w(x)>0$, followed by $w(x)\le1$, proves
\eqref{ct:localized-initial}.  Integrate
\eqref{ct:localized-source}, use the cumulative contact inequality and
the nonnegativity of the corresponding integral of $h$, and again
divide by the fixed number $w(x)$.  This proves
\eqref{ct:localized-cumulative}.
\end{proof}

The elementary proof just given requires no positive lower bound on
$s$.  In applications of first variation we also cut off the Gaussian
test field within the original chart.  If $s\le\rho\ll R$ and its
center lies in the fixed interior region, this introduces a factor
$\exp(-cR^2/s^2)$ times the bounded chart mass and fixed powers of
$s^{-1}$.  The same reasoning makes these errors smaller than any
prescribed part of the subtractive margin.  At a node of the recursion
below, $s\ge r_b$ and $\rho/R=T^{-4N}$, so this remains true even after
the factor $s^{-m}\le\exp(O(T))$ is included.  All subsequent Gaussian
identities use this convention.

\subsection{The estimate for separated configurations}

We use the signed estimate \Cref{sx:signed} only at smaller integers in
this subsection.  Separation refers to the scaled normal coordinates;
the integrands defining $\Lambda$ continue to use physical units.
In this subsection write $z=\zeta=z_{\rm phys}/l$, in the root chart;
these height coordinates remain fixed at all recursive nodes.

For two lists $\mathcal A=\{a_1,\ldots,a_Q\}$ and
$\mathcal B=\{b_1,\ldots,b_Q\}$, write
\[
 d_\infty(\mathcal A,\mathcal B)
 =\min_{\sigma\in\mathfrak S_Q}\max_{1\le j\le Q}|a_j-b_{\sigma(j)}|.
\]
Thus repetitions are distinct entries in a matching. Choose once
$0<r_{\rm maj}<1$ so that the fixed enlargements needed first to select
the integer by \Cref{lt:scalar} and then to apply \Cref{lt:config} on
the inner ball $B_{r_{\rm maj}}$ lie in $B_4$.
The corresponding inner ball of a work ball $B_r(y)$ is
$B_{r_{\rm maj}r}(y)$.

We use comparison chains whose work radii change by factors at most two
and whose center steps are at most one quarter of the smaller inner
radius. For adjacent inner balls $B,B'$, these choices give a fixed
$\vartheta_{\rm ov}>0$ such that
\begin{equation}\label{ct:majority-overlap}
 \vol_S(B\cap B')\ge
 \vartheta_{\rm ov}\bigl(\vol_S(B)+\vol_S(B')\bigr).
\end{equation}
Indeed their intersection contains a ball of a fixed fraction of the
smaller radius, and their radii are comparable. Fix
$0<\delta_{\rm maj}<\vartheta_{\rm ov}$, and let $C_{\rm maj}$ be a fixed
constant valid in \Cref{lt:config} for this exceptional proportion and
the finitely many fixed enlargements used below. Every auxiliary
configuration below uses this same exceptional proportion. The corresponding
smallness threshold in \Cref{lt:config} is therefore fixed before $T$;
the local bounds used to select these configurations tend to zero with $T$.

Suppose $s^{-m}\tau(B_{cs}(x))\le d_{\rm loc}$ on a ball and its required
fixed enlargement. A \emph{majority configuration} on $B_s(x)$ is a list
of $Q$ heights such that, outside a set of base volume at most
$\delta_{\rm maj}\vol_S(B_s(x))$, each fiber has exactly $Q$ projection
entries and these entries match the list within
$C_{\rm maj}s\sqrt{d_{\rm loc}}$. The projection entries repeat each
spatial point according to its integer multiplicity. If majority
configurations $\mathcal A,\mathcal B$ on adjacent balls of radii $s,t$
have local upper bounds $d_s,d_t$, then
\begin{equation}\label{ct:majority-matching}
 d_\infty(\mathcal A,\mathcal B)
 \le C_{\rm maj}\bigl(s\sqrt{d_s}+t\sqrt{d_t}\bigr).
\end{equation}
In fact, \eqref{ct:majority-overlap} leaves a fiber outside both
exceptional sets. Composing its two matchings proves the estimate,
including repetitions.

At the root the inner ball is $B_{r_{\rm maj}}$ and
$d_{\rm loc}=r_{\rm maj}^{-m}d$, so its matching error is a fixed
multiple of $\sqrt d$. The majority property itself imposes no condition
on exceptional fibers; their contribution remains in the global $\tau$
budget. A partition of a list's entries into groups is separated by
$\delta$ if different groups have mutual distance at least $\delta$.
Identical list entries therefore belong to the same group when $\delta>0$.

\begin{proposition}[Separated configurations]\label{ct:separated}
Fix $Q>1$ and assume \Cref{sx:signed} for every integer $1\le q<Q$.
For each fixed $P>0$ there are constants $C_Q,C_S$ and a function
$b=b_Q$ satisfying
\begin{equation}\label{ct:b-choice}
 b(T)\gg(\log T)^{200},\qquad b(T)=T^{o(1)},\qquad
 F_{Q-1}(b(T)/2)\ge(\log T)^3
\end{equation}
such that the following holds.  Under the hypotheses of
\Cref{sx:signed}, let $\mathcal A=\{a_1,\ldots,a_Q\}$ be any majority
configuration on $B_{r_{\rm maj}}$ with local upper bound
$d_{\rm loc}=r_{\rm maj}^{-m}d$, and suppose it has a partition into at least two
nonempty groups separated by
$10S(T)\sqrt d$, where
\[
 S(T)=e^{C_Sb(T)},\qquad c(T)=\frac{(\log T)^2}{T}.
\]
Let $D_{\rm far}$ be the projection of tilt restricted to
$\{\dist(z,\mathcal A)>S(T)\sqrt d\}$.  Then every nonnegative weight
$w\le1$, supported in $B_1$ with $\Lip w\le1$, satisfies
\begin{equation}\label{ct:separated-estimate}
 \int w\,\dd\bigl(L-C_Q\Lambda+c(T)D_{\rm far}\bigr)
 \le T^{-P}d.
\end{equation}
The statement holds on other fixed interior balls after rescaling.
Polynomial factors in the Lipschitz norm or the enlargement of a
Gaussian working ball can be allowed by increasing $P,C_S$ and the
fixed constants.
\end{proposition}

\begin{proof}
The proof is a finite recursion over working balls, with the root
configuration and its far-tilt region kept fixed. On a sufficiently
small ball, the support lies in disjoint neighborhoods of the root
groups. Cutting through the gaps gives stationary pieces with positive
counts below $Q$, so the lower-integer theorem makes the signed integral
small on that ball.

On a larger ball, if the signed integral exceeds a small margin, the
remaining positive source gives a contact measure.
The cumulative contact inequality and the Gaussian mass identity exclude
a contact whose derivative scale is close to a threshold hit or a side
equality: it would force too much signed mass at the outer scale.
Bottom contacts are excluded by the small-ball reduction. Each remaining
contact lies on a vertical side, far below its equality scale. Its signed
Gaussian integral is partitioned over smaller working balls, while a choice of thresholds
made before constructing the contact measure confines the associated
charging balls to a small part of the measure on the working ball. Thus a large
generation pays a small margin and at most half of the next-generation
bound. The choices below enforce these alternatives and keep the
configuration motion smaller than the original separation.

For example, after extending $F_{Q-1}$ monotonically over a fixed
initial interval, one can use
\[
 b(T)=1+(\log T)^{201}
       +4F_{Q-1}^{-1}((\log T)^3).
\]
The definition of the iterated exponentials gives
$F_{Q-1}^{-1}((\log T)^3)=T^{o(1)}$, and this choice has
\eqref{ct:b-choice}.  All changes of $b$ below are by fixed constants
or an additional fixed power of $\log T$.

Choose a fixed integer $N>P+10$, increasing it later if necessary, and
then choose a fixed $J_0$ sufficiently large in terms of $m,n,Q,N$.
In particular require
\begin{equation}\label{ct:J-choice}
 J_0>4N(m+2)+100,
\end{equation}
and enlarge $J_0$ to exceed all fixed polynomial density losses in the
cell constructions below.  Set
\begin{equation}\label{ct:recursion-parameters}
 \eta=T^{-N},\quad \varepsilon=\frac{(\log T)^{60}}{T},\quad
 r_b=\sqrt d\,e^{3b},\quad n_0=\lceil10N\log T\rceil,
 \quad E_j(t)=T^{J_0(j+1)}d\,t^{-2-\varepsilon}.
\end{equation}
The parameter $T$, the original configuration $\mathcal A$, and the
measure $D_{\rm far}$ remain fixed throughout the finite recursion.
They are never recomputed at a child node.

Two elementary parameter bounds will be used repeatedly.  With
$\ell_T=\log T$, direct substitution gives
\begin{equation}\label{ct:small-envelope}
 \log E_{2n_0+4}(r_b/T^4)
 =-6b+\tfrac12\ell_T^{60}-3\varepsilon b
       +O_{J_0,N}(\ell_T^2)
 <-2b.
\end{equation}
Also, for $i\le n_0$,
\begin{equation}\label{ct:motion-budget}
 v\sqrt{E_{2i+1}(v)}
 =\sqrt d\,T^{J_0(i+1)}v^{-\varepsilon/2}.
\end{equation}
For $v\ge cr_b/T^2$, the logarithm of the factor multiplying $\sqrt d$
is $O_{J_0,N}((\log T)^2)+O((\log T)^{60})$.  Consequently $O(T)$
such increments, even with any fixed polynomial loss in $T$, sum to
less than $e^b\sqrt d$.  This is why the much larger separation
$e^{C_Sb}\sqrt d$ survives every generation.

At generation $i$ a node consists of a ball $B_{4R}(x_*)$ with
\begin{equation}\label{ct:node-hypotheses}
 \tau(B_{4R}(x_*))\le C R^mE_{2i}(R),\qquad
 R\ge cr_b/T^2,
\end{equation}
and a majority configuration $\mathcal A_R$ on
$B_{r_{\rm maj}R}(x_*)$, with local upper bound
$C r_{\rm maj}^{-m}E_{2i}(R)$, satisfying
\[
 d_\infty(\mathcal A_R,\mathcal A)\le(i+1)e^b\sqrt d.
\]
The integer is $Q$ on the balls constructing the node; the small-density
and overlap checks below establish this before their configurations are
selected. We estimate all weights
$0\le w\le1$, $\Lip w\le R^{-1}$, supported in $B_R(x_*)$.
Write
\[
 L'=L-C_Q\Lambda+c(T)D_{\rm far},\qquad
 \mathcal N_R=\tau(B_{4R}(x_*))+dR^m.
\]
We prove, by downward induction on $i$, that
\begin{equation}\label{ct:node-conclusion}
 \int w\,\dd L'\le p_i\mathcal N_R,
 \qquad p_{n_0}=C_*,\qquad
 p_i=C\eta+\tfrac12p_{i+1}.
\end{equation}
Here $C_*,C$ are fixed, independent of $i$ and $T$.  The last-generation
bound follows at once from $|L'|\le C\tau$.  Initialization at $R=1$
follows from $\tau(B_4)\le d$ and the given configuration.

\paragraph{Small nodes and actual stationary pieces.}
Suppose first that $R\le T^{10N}r_b$, and put $E=E_{2i}(R)<1$,
as follows from \eqref{ct:small-envelope}.
Cover $B_{2R}(x_*)$ by a fixed finite number of balls of radius
$\sigma R$, including the concentric ball, where the fixed $\sigma>0$
is small enough that every
fixed enlargement needed for scalar selection and then for the support
estimate in \Cref{lt:config} lies in $B_{4R}(x_*)$, and the scalar
enlargement of the concentric ball lies in $B_{2R}(x_*)$.
Their normalized $\tau$ is at most $CE$.
The normalized smooth mass averages therefore differ from $Q$ by
$O(E)<1/4$, so \Cref{lt:scalar} selects $Q$ on these balls.
On each choose the joint configuration in \Cref{lt:config} with
exceptional proportion $\delta_{\rm maj}$, retaining its majority property
and using its all-support bound with a fixed
$0<\beta<\min\{1/2,1/m\}$. Its local upper bound is $CE=o(1)$.
Connect it to $\mathcal A_R$ by a fixed chain satisfying
\eqref{ct:majority-overlap}. The intermediate balls have the same
$CE=o(1)$ upper bound, and the same smooth mass test selects $Q$ there.
Equation \eqref{ct:majority-matching}
bounds this change by $CR\sqrt E$, while the support bound is $CRE^\beta$.
Consequently all support heights over $B_{2R}(x_*)$ satisfy
\begin{equation}\label{ct:bottom-sup}
 \begin{split}
 \dist(z,\mathcal A)
 &\le (i+1)e^b\sqrt d+CR\sqrt E+CRE^\beta\\
 &\le (i+1)e^b\sqrt d+C R E_{2i}(R)^\beta
 \le e^{C'b}\sqrt d
 \end{split}
\end{equation}
for a fixed $C'$. The second line uses $\beta<1/2$ and $E<1$.
To see the last bound without a
hidden scale restriction, write $R=\sqrt d\,e^{3b}T^\theta$, where
$-3\le\theta\le10N$, and substitute in $RE_{2i}(R)^\beta$;
its logarithm after division by $\sqrt d$ is bounded above by $C'b$
using \eqref{ct:small-envelope}.  Choose $C_S>C'+10$.
For a root group $\mathcal A_g$, put
\[
 \mathcal O_g=\{z:\dist(z,\mathcal A_g)<2S(T)\sqrt d\}.
\]
For large $T$, \eqref{ct:bottom-sup} puts all the support over $B_{2R}$
inside the smaller regions with radius $S(T)\sqrt d$. The regions
$\mathcal O_g$ are disjoint, since the root group gap is
$10S(T)\sqrt d$. Thus there is an open height gap around each group's
support, and no far energy in this interior region.

The use of disjoint height regions to obtain stationary restrictions
has a precedent in \cite[Section~3.1]{BDD2025}. For each group, restrict
the varifold over $B_{2R}(x_*)$ to $\{z\in\mathcal O_g\}$. Choose a
smooth height cutoff equal to one near the smaller region of radius
$S(T)\sqrt d$ around $\mathcal A_g$ and zero outside $\mathcal O_g$.
Its derivative vanishes near the entire support over $B_{2R}(x_*)$.
Multiplying any ambient test field with compact base support there by
this cutoff proves stationarity of the restriction. It keeps every
included spatial point with
its entire integer multiplicity, including on exceptional fibers.
On a majority fiber of one of the local configurations, its matching
to the fixed root list moves each entry by at most
$(i+1)e^b\sqrt d+CR\sqrt E=o(S(T)\sqrt d)$. An entry matched to
$\mathcal A_g$ therefore lies in $\mathcal O_g$. If two copies of one
spatial point were matched to different root groups, their two root
entries would be separated by at most twice this error, less than
$10S(T)\sqrt d$ for large $T$, contrary to the separation.
Each group's projection count on that fiber is thus
the number $q_g$ of entries in its root group, including repetitions.
Since there are at least two nonempty groups,
$1\le q_g<Q$ and $\sum_gq_g=Q$. To specify the individual energy budgets,
take a finite nonnegative partition of the weight on balls
$B_{r_\alpha}(y_\alpha)$, with $r_\alpha=cR$ for a sufficiently small
fixed $c$, and $B_{8r_\alpha}(y_\alpha)\subset B_{2R}(x_*)$.
These enlarged balls have bounded overlap. For each stationary
restriction, its normalized projected mass is bounded by the normalized total
mass, and its normalized $D+\Lambda$ is $O(E)=o(1)$. Apply
\Cref{lt:scalar} with its initial intermediate ball equal to the
concentric majority ball above. Its selected integer is $q_g$ by
intersection with the group majority set there. Connect this ball to
the intermediate balls for the scalar applications on
$B_{4r_\alpha}(y_\alpha)$ by a fixed finite chain of comparable
interior balls with the positive overlap in \eqref{ct:majority-overlap}.
The scalar count exceptions on each intermediate ball have relative
volume $O(E)$, smaller than the fixed overlap for
large $T$; on a remaining fiber the two selected integers coincide.
Thus $q_g$ is identified and propagated before the mass defect of the
restriction is estimated.

Write $\mu_g$ for the mass measure of the restriction and $D_g,\Lambda_g$
for its projected energy measures. Now set
$M_g=\pi_\#\mu_g-q_g\vol_S$ and $\tau_g=|M_g|+D_g+\Lambda_g$.
The scalar estimate with interior
room now gives
\[
 \tau_g(B_{4r_\alpha}(y_\alpha))
 \le C_s(D_g+\Lambda_g)(B_{8r_\alpha}(y_\alpha)).
\]
Define the parameter for the lower-integer application of \Cref{sx:signed} by
\begin{equation}\label{ct:group-budget}
 d_{g\alpha}=\max\left\{d,
 C_s r_\alpha^{-m}(D_g+\Lambda_g)
                  (B_{8r_\alpha}(y_\alpha))\right\},
 \qquad T_{g\alpha}=-\log d_{g\alpha}.
\end{equation}
Increasing the fixed $C_s$ if necessary supplies the required bound
for $\tau_g$ on each rescaled $B_4$.  Moreover
\begin{equation}\label{ct:group-budget-sum}
 \sum_{g,\alpha}d_{g\alpha}r_\alpha^m
 \le C\bigl((D+\Lambda)(B_{4R})+dR^m\bigr)
 \le C\mathcal N_R.
\end{equation}
Here the group energies add exactly on their separated region, and
the enlarged base balls have bounded overlap.  Individual mass defects
have been estimated by \Cref{lt:scalar}, rather than by splitting
$|M|$.  The parameters satisfy
\begin{equation}\label{ct:child-parameter}
 b/2\le T_{g\alpha}\le T.
\end{equation}
Indeed $d\le d_{g\alpha}\le d+C E_{2i}(R)$, and
\eqref{ct:small-envelope} absorbs the fixed constants.  Physical
radii only decrease, and physical heights are unchanged, so the
curvature hypothesis follows from the original one.  Apply
\Cref{sx:signed} at $q_g$ to the partition weights normalized by a
fixed constant.
Taking $C_Q$ at
least as large as the finitely many constants for $q<Q$, summing the
lower-integer conclusions with \eqref{ct:group-budget-sum} gives an
$o(\eta)\mathcal N_R$ bound.  Indeed
$F_q(T_{g\alpha})\ge F_{Q-1}(b/2)\ge(\log T)^3$ for large $T$.
This argument also applies to the small cells used at a bottom contact.

\paragraph{Stopping tests at a large node.}
Now assume $R>T^{10N}r_b$.  Use a torus of period $C R$, identifying
the node and a fixed surrounding chart with its central portion.  Put
\[
 \nu=\tau|_{B_{4R}(x_*)}+d\,\dd y,
 \qquad \rho=T^{-4N}R.
\]
Its total mass is comparable to $\mathcal N_R$.  If
$f=wL'-\eta\nu$ has nonpositive integral there is nothing to prove.
Otherwise apply \Cref{ct:contact} with bottom $r_b$ and dyadic test
scales $t\in[r_b,\rho]$, using
\begin{equation}\label{ct:stopping-q}
 q_t(x)=\frac{(\tau|_{B_{4R}(x_*)})_t(x)}{E_{2i+1}(t)},
 \qquad q_t(x)\le\lambda_t,\quad 1\le\lambda_t\le2.
\end{equation}
The thresholds will be chosen later.  There are $O(T)$ test scales.
The torus masses satisfy the growth condition in
\Cref{ct:localization}, because
$\mathcal N_R/(dR^m)\le\exp(O(T))$ by
\eqref{ct:node-hypotheses}.  Also
$\rho/R=T^{-4N}\ll\eta^2/T^2$.

For every contact $(x,s_0)$, localization places $x$ in the interior
support region.  Gaussian comparison between successive dyadic scales,
and the non-strict tests at larger scales, give
\begin{equation}\label{ct:contact-envelope}
 \tau_v(x)\le C E_{2i+1}(v)\qquad(s_0\le v\le\rho).
\end{equation}
Here and below the measure in this display is the node restriction.
At the current mesh scale its value is bounded by the next larger
scale.  At the top use the node's total mass.  In fact, for
$v\ge\rho/16$,
\[
 q_v(x)\le C T^{-J_0}(R/v)^{m-2-\varepsilon}
 \le C T^{-J_0+4N\max(m-2,0)}.
\]
Thus all top tests hold strictly, and every scale with $q_v\ge1$
lies below $\rho/16$, leaving a fixed outward interval below $\rho$.

If $s_0>r_b$, the boundary description
\eqref{ct:mesh-domain} supplies a mesh scale $v_1\ge s_0$ with
$q_{v_1}(x)\ge1$.  Either it is a current hit, or it is a larger scale
with $q_{v_1}(x)=\lambda_{v_1}$.  In the latter case every still larger
test satisfies its continuation inequality.  This follows directly
from the finite intersection: away from the bottom, if there were no
current obstruction and every relevant larger test were strict, a
space--scale neighborhood would belong to $\mathcal D$.

\paragraph{Exclusion of a nearby hit or side equality.}
We show that no such $v_1$ can satisfy
$v_1/s_0\le(\log T)^{20}$.  All the following estimates are on this
one outward vertical.  Put
\[
 I_1=v_1^2E_{2i+1}(v_1),\qquad A(v)=v^2M_v(x),\qquad
 Z(s)=\int_{s_0}^s v^2(L'-\tfrac\eta4\nu)_v(x)\,\dd\log v.
\]
By \Cref{ct:localization}, $Z\ge0$ and $A(s_0)\ge0$; the latter
uses $c(T)<1$ and $L'=2M-D-C_Q\Lambda+c(T)D_{\rm far}$.
The envelope gives
\begin{equation}\label{ct:nearhit-cost}
 \int_{s_0}^{\rho}v^2D_v\,\dd\log v
 \le \frac{C I_1}{\varepsilon}
          (v_1/s_0)^\varepsilon
 \le \frac{C I_1}{\varepsilon},
 \qquad \frac{c(T)}\varepsilon=(\log T)^{-58}.
\end{equation}
In particular the adverse $c(T)D_{\rm far}$ term is small enough for
the logarithmic losses below.

At $v_1$, and on its fixed outward dilations, the scalar estimate is
available in weighted form:
\begin{equation}\label{ct:nearhit-scalar}
 |M|_{v_1}\le C(D+\Lambda)_{v_1}+o(E_{2i+1}(v_1)).
\end{equation}
Here are the raw density and integer checks, also needed for the
children below.  For $s\in[s_0,\rho]$ and
$U=1+|y-x|/s\le C\sqrt T$, put $r=cs/U$ with a small fixed $c$.
All required fixed enlargements $B_{Lr}(y)$ lie in $B_{2R}(x_*)$:
$x\in B_R(x_*)$ and $C\sqrt T\,\rho=o(R)$.  Choose
$a=C_1Us$ so that $G_a(x-z)\ge c_1a^{-m}$ on these enlarged balls.
If $a\le\rho$, positivity and the envelope give
\[
 r^{-m}\tau(B_{Lr}(y))
 \le C(a/r)^m E_{2i+1}(a)
 \le C U^{2m-2-\varepsilon}E_{2i+1}(s).
\]
If $a>\rho$, use the node total instead.  Dividing by
$E_{2i+1}(s)$ gives the bound
\[
 C T^{-J_0}U^m(R/s)^{m-2-\varepsilon}
 \le C T^{-J_0+4N\max(m-2,0)}
             U^{m+\max(m-2,0)}\le C.
\]
In the middle inequality a negative exponent of $R/s$ is discarded;
for a positive exponent use $R/s\le C U T^{4N}$.
Consequently, in both cases,
\begin{equation}\label{ct:cell-density}
 r^{-m}\tau(B_{Lr}(y))
 \le C T^{\max(m-1,0)}E_{2i+1}(s)=o(1).
\end{equation}
Since $\tau$ includes $|M|=|\pi_\#\mu-Q\vol_S|$, choose a fixed
nonnegative smooth mass test on this enlarged ball, normalized to
reference integral one and with bounded scaled derivatives.  Its
average differs from $Q$ by at most the right side times a fixed
constant, hence by less than $1/4$.  This identifies $Q$ directly
in \Cref{lt:scalar}, and also
provides its mass and energy hypotheses.  Apply \Cref{lt:scalar} at $s=v_1$
on a cover by these cells and sum with the comparable Gaussian
weights, as in \Cref{lt:weightedballs}.  The cutoff at
$C\sqrt T\,v_1$ has arbitrarily small exponential tail for a large
fixed $C$, using the node total and $v_1\ge r_b$.

The identity in \Cref{lt:gaussian}, integrated using the nonnegative
bank $Z$, gives
\begin{equation}\label{ct:nearhit-bank}
 A(s)\ge A(s_0)+\int_{s_0}^s v^2R_v\,\dd\log v+Z(s)
             -C(\log T)^{-58}I_1,
 \qquad s_0\le s\le\rho.
\end{equation}
The geometric errors are absorbed by part of the remaining margin
$\eta\nu/4$: on $|u|\le T$ they carry the factor
$Ce^{-\kappa T}(1+T^2)=o(\eta)$, and the tails are negligible under
\eqref{ct:contact-envelope} and the node total.  The nonnegative
$C_Q\Lambda$ term has been discarded.

There is a fixed $c_1>0$ such that $A(v_1)\ge c_1I_1$.
This includes a current hit with $q_{v_1}>\lambda_{v_1}$: although
its current test is omitted, the next larger mesh scale $t_+\le2v_1$
obeys $\tau_{t_+}\le2E_{2i+1}(t_+)$, and
$G_{v_1}\le(t_+/v_1)^mG_{t_+}$ bounds its current average.
Only $q_{v_1}\ge1$ is needed for the lower bound.

For an explicit cumulative-sign argument, fix $C_Q\ge1$ first and
let $C_s$ be the constant in \eqref{ct:nearhit-scalar}.  Write
$E_1=E_{2i+1}(v_1)$ and take $T$ large enough that its tail is at
most $E_1/2$ and $c(T)\le1/2$.  Then
\[
 (D+\Lambda)_{v_1}\ge\alpha E_1,
 \qquad \alpha=\frac{1}{2(C_s+1)}.
\]
If $A(v_1)\le c_1I_1$ with $c_1\le\alpha/8$, then
$L'_{v_1}\le2c_1E_1-(D+\Lambda)_{v_1}/2\le-\beta E_1$,
where $\beta=\alpha/4$.  For
$v\in[v_1,2^{1/4}v_1]$, differentiating the fixed source $L'$
and dominating $|\dot G_v|$ by $CG_{2v}$ gives
$|\dot L'_v|\le K_1E_1$.  The constant $K_1$ is fixed after $C_Q$;
the envelope at $2v$ follows from strictly larger mesh tests, and
$v_1<\rho/16$ leaves room.  In particular this estimate does not
use the omitted current test or differentiate the fixed far cutoff.
Choose
\[
 \theta=\min\{\tfrac14\log2,\beta/(2K_1)\},\qquad
 \gamma=\beta\theta/2,\qquad
 0<c_1\le\min\{\alpha/8,\gamma/4\}.
\]
On $[v_1,e^\theta v_1]$ one has $L'_v\le-\beta E_1/2$,
hence $\dot Z\le-\beta I_1/2$.  Nonnegativity of $Z$ at the
right endpoint gives $Z(v_1)\ge\gamma I_1$.
Finally take $T$ large enough that the error in
\eqref{ct:nearhit-bank} is at most $\gamma I_1/2$.
That inequality at $v_1$ yields $A(v_1)\ge\gamma I_1/2>c_1I_1$,
a contradiction.  All constants were fixed before this last largeness
threshold; only the cumulative sign of $Z$ was used at later scales.

The upper differential inequality
$\dot A\le2A+v^2R_v+$ absorbed errors, again from
\Cref{lt:gaussian}, and $v_1/s_0\le(\log T)^{20}$ now yield
\begin{equation}\label{ct:nearhit-surplus}
 A(s_0)+\int_{s_0}^{v_1}v^2R_v\,\dd\log v
 \ge c(\log T)^{-40}I_1.
\end{equation}
One may retain the geometric error explicitly in the integrating
factor formula: divided by $I_1$ it is
$O(e^{-\kappa T}T^C/\varepsilon)$, hence smaller than every fixed
inverse power of $\log T$.  Inserting
\eqref{ct:nearhit-surplus} in \eqref{ct:nearhit-bank} at $\rho$
gives $A(\rho)\ge c(\log T)^{-40}I_1$.  On the other hand the
node mass gives
\[
 A(\rho)\le C\rho^2(R/\rho)^m E_{2i}(R)
   = C T^{J_0(2i+1)+4N(m-2)}dR^{-\varepsilon},
 \qquad I_1=T^{J_0(2i+2)}dv_1^{-\varepsilon}.
\]
The ratio of the upper bound to $I_1$ is at most
$CT^{-J_0+4N(m-2)}$, contradicting \eqref{ct:J-choice}.
This excludes current hits and all side equalities within the stated
logarithmic factor.

\paragraph{Gaussian children and the bottom.}
For any remaining contact, configurations on the Gaussian cells at
$s_0$ match the node configuration with movement at most $e^b\sqrt d$.
Here is the comparison chain for a cell with work radius comparable to
$s_0/U$, where $U=1+|y-x|/s_0\le C\sqrt T$. First shrink concentrically
at $x_*$ to a small fixed multiple of $R$, move the center to $x$ at
that radius, and then shrink concentrically to a work radius comparable
to $Us_0$. Move the center to $y$ at that radius and finally shrink to
the cell radius. At each stage change radii by factors at most two and
make center steps at most one quarter of the smaller inner majority
radius. All these inner balls satisfy \eqref{ct:majority-overlap}, so
\eqref{ct:majority-matching} bounds the changes by the sum of their
local matching errors. The initial fixed chain stays inside
$B_{4R}(x_*)$; the off-center chain at scale $s_0$ lies in
$B_{CUs_0}(x)\subset B_{2R}(x_*)$. The envelope bounds these increments by
\eqref{ct:motion-budget}, with fixed polynomial losses.  More
explicitly, above $\rho$ the node total gives, for
$\rho\le t\le cR$,
\[
 t\sqrt{t^{-m}\tau(B_{ct}(x))}
 \le C R\sqrt{E_{2i}(R)}\,
              T^{2N\max(m-2,0)}.
\]
Off-center balls nested from scale $Us_0$ down to $s_0/U$ have the
same increment bound as $s_0\sqrt{E_{2i+1}(s_0)}$, with another
fixed power of $U\le C\sqrt T$, by the raw estimates proving
\eqref{ct:cell-density}. All normalized bounds used for intermediate
configurations are at most a fixed power of $T$ times $E_{2i}(R)$ or
$E_{2i+1}(t)$ for $t\ge cr_b/T^2$. Hence they tend to zero by
\eqref{ct:small-envelope}. Their $|M|$ terms give the same smooth mass
test selecting $Q$, and \Cref{lt:config} is applied with the fixed
$\delta_{\rm maj}$ before each matching. The logarithms of all these factors are
$O_{J_0,N,m}((\log T)^2)+O((\log T)^{60})$ after division by
$\sqrt d$.  Since there are $O(T)$ steps,
\eqref{ct:motion-budget} proves the claimed $e^b\sqrt d$ allowance.

Here is the precise normalization of the cells.  A smooth partition of
unity on $|u|\le C\sqrt T$, where $u=(y-x)/s_0$, can be chosen with
bounded overlap on balls $B_{r_j}(y_j)$, with
\begin{equation}\label{ct:gaussian-cells}
 r_j\asymp\frac{s_0}{1+|y_j-x|/s_0},\qquad
 |\nabla\chi_j|\le C/r_j.
\end{equation}
Use a sufficiently small fixed constant in $\asymp$ to leave room for
$B_{4r_j}(y_j)$.  This follows, for example, by a maximal disjoint
selection for the Lipschitz radius function in
\eqref{ct:gaussian-cells}, followed by normalized smooth bumps.
On the enlarged ball, $G_{s_0}(x-y)$ differs from
$G_{s_0}(x-y_j)$ by a fixed factor, and its logarithmic gradient is
$O(r_j^{-1})$.  Thus
\[
 w_j(y)=\frac{\chi_j(y)G_{s_0}(x-y)}{C G_{s_0}(x-y_j)}
\]
is an admissible child weight after a fixed normalization.  In
particular, this normalization has no power of $T$.

The child's unweighted density was estimated in
\eqref{ct:cell-density}.  The identity
\[
 E_{2i+2}(r_j)=T^{J_0}E_{2i+1}(s_0)(s_0/r_j)^{2+\varepsilon}
\]
then gives
\[
 r_j^{-m}\tau(B_{4r_j}(y_j))
 \le C T^{\max(m-1,0)}E_{2i+1}(s_0)
 \le C E_{2i+2}(r_j),
\]
by the choice of $J_0$, with $C$ independent of $i,T$.
Also $r_j\ge cr_b/\sqrt T\ge cr_b/T^2$.
The smooth mass test following \eqref{ct:cell-density} identifies $Q$
on each child. After scaling its work ball to one, choose from
\Cref{lt:config} a majority configuration on $B_{r_{\rm maj}}$
with local upper bound $C r_{\rm maj}^{-m}E_{2i+2}(r_j)$.
Using the same normal frame, multiply its heights by $r_j$ to return
to the root coordinates.
The matching argument above gives its allowance relative to
$\mathcal A$, so these are admissible next-generation
nodes.  Although the raw density estimate loses a polynomial factor,
summation of the resulting estimates does not: bounded overlap and
Gaussian comparability give
\begin{equation}\label{ct:cell-sum}
 \sum_jG_{s_0}(x-y_j)
       \bigl[\tau(B_{4r_j}(y_j))+dr_j^m\bigr]
 \le C\bigl(\tau_{s_0}(x)+d\bigr).
\end{equation}
The fixed $d$ floors therefore sum to the correct constant-density
Gaussian average.

At a bottom contact $s_0=r_b$, these cells are small nodes.  The
stationary-piece argument above, including
\eqref{ct:child-parameter}, bounds $L'_{s_0}$ by
$o(\eta)\nu_{s_0}$.  The Gaussian tail outside $C\sqrt T$ is
$o(\eta)d$ for a sufficiently large fixed $C$: combine its
$e^{-C^2T/4}$ factor with total node mass and
$s_0^{-m}\le\exp(O(T))$.  This contradicts
\eqref{ct:localized-initial}; there are no bottom contacts.
At any other remaining contact the next-generation bound and
\eqref{ct:cell-sum} give
\begin{equation}\label{ct:child-contact-bound}
 h_{s_0}(x)\le C_0p_{i+1}\nu_{s_0}(x).
\end{equation}
All omitted tails are absorbed by the subtractive margin in
\eqref{ct:localized-initial}.  In particular, every remaining contact
is a vertical-side contact with some equality scale
\begin{equation}\label{ct:remote-side}
 v_1/s_0>(\log T)^{20}.
\end{equation}

\paragraph{First charging step: replace a Gaussian by a nearby ball.}
Put $J=\log T$ and $p=p_{i+1}$.  Call a contact $(x,s)$ good if some
closed ball of radius $0<r\le Js$ satisfies
\begin{equation}\label{ct:good-ball}
 h(\overline B_r(x))\le A_0p\nu(\overline B_r(x)),
\end{equation}
where $A_0\ge2C_0$ is a fixed constant.  Otherwise call it bad.
We first show that all bad centers together carry only
$o(\eta)\mathcal N_R$ mass of $h$.

For the Euclidean Gaussian, radial layer cake gives
\[
 \nu_s(x)=c_ms^{-m}\int_0^\infty
                  t e^{-t^2/2}\nu(B_{ts}(x))\,\dd t.
\]
At a bad contact the failure of \eqref{ct:good-ball} for every
$r\le Js$ bounds the part with $t\le J$ by $h_s/(A_0p)$.
Define the remaining layer-cake tail by
\[
 \mathcal T_{J,s}\nu(x)=c_ms^{-m}\int_J^\infty
                   t e^{-t^2/2}\nu(B_{ts}(x))\,\dd t.
\]
This is a tail in the layer-cake radius: it includes the residual
weight $c_ms^{-m}e^{-J^2/2}\nu(B_{Js}(x))$ as well as mass farther
away.  Combining
$\nu_s\le h_s/(A_0p)+\mathcal T_{J,s}\nu$ with
\eqref{ct:child-contact-bound} gives
\begin{equation}\label{ct:bad-tail}
 h(B_{5s}(x))\le C_m s^m h_s(x)
              \le C_mp s^m\mathcal T_{J,s}\nu(x).
\end{equation}

Let $E_{\rm bad}$ consist of spatial contact centers that admit no
good representative; for each of them assign one contact scale.
The elementary $5r$ covering lemma selects disjoint assigned balls
whose fivefold enlargements cover these centers.  Their scales lie in
$[r_b,\rho]$, hence in $O(T)$ dyadic bins.  In a bin $[a,2a)$,
disjointness implies, for every $y$ and $t\ge1$,
\[
 \sum_j\mathbf1_{B_{ts_j}(x_j)}(y)\le C_m(1+t)^m:
\]
the disjoint balls of radius at least $a$ have centers in a ball of
radius $Cta$.  Sum \eqref{ct:bad-tail}, use Tonelli, and then sum the
bins to obtain
\begin{equation}\label{ct:bad-total}
 h^*(E_{\rm bad})
 \le C_mpT(1+J)^{m+1}e^{-J^2/2}\nu(X)
 =o(\eta)\mathcal N_R.
\end{equation}
Here $h^*$ denotes outer measure, so no measurable choice of a contact
above each center is needed.  The same proof on the torus uses periodic
lifts of $h,\nu$: the small selected balls have disjoint lifts, the
packing estimate is unchanged, and the periodic Gaussian is exactly
the sum over those lifts.  Alternatively the nonnearest copies can be
absorbed in the exponentially smaller chart tail.

\paragraph{Second charging step: localize the union of good balls.}
Every good ball lies in a set determined only by the node's tests and
thresholds.  To see this, let $(x,s)$ be its contact and choose the
equality scale $v_1$ in \eqref{ct:remote-side}.  For every mesh scale
$v\ge v_1$ and every $|y-x|\le Js$, Gaussian differentiation gives
\begin{equation}\label{ct:q-gradient}
 |\nabla q_v(y)|\le C/v.
\end{equation}
Indeed $|\nabla G_v(y-z)|$ is bounded by
$Cv^{-1}G_{2v}(x-z)$ when $|y-x|\le v/10$; the envelope at $2v$
proves the assertion when $2v\le\rho$.  For $2v>\rho$ the node total
and the top slack following \eqref{ct:stopping-q} give the same bound.
Since $Js/v_1<J^{-19}$, for sufficiently large $T$ every point of
the good ball satisfies
\begin{equation}\label{ct:charged-tests}
 q_v(y)\le\lambda_v+\delta\quad(v>v_1),\qquad
 |q_{v_1}(y)-\lambda_{v_1}|\le\delta,
 \qquad\delta=J^{-15}.
\end{equation}

Order the mesh scales $v_1',\ldots,v_K'$ from largest to smallest,
and define the Borel set
\begin{equation}\label{ct:threshold-event}
 \mathcal E_\lambda=
 \bigcup_{k=1}^K\left[
  \{|q_{v_k'}-\lambda_{v_k'}|\le\delta\}
  \cap\bigcap_{j<k}\{q_{v_j'}\le\lambda_{v_j'}+\delta\}
                    \right].
\end{equation}
Equation \eqref{ct:charged-tests} says that every original good ball,
not just its center, lies in $\mathcal E_\lambda$.

We use the following form of the Besicovitch covering theorem
\cite[Chapter~1, Lemma~4.5]{Simon2018}.
For a bounded set with one assigned positive-radius ball at each
point and bounded assigned radii, there is a countable subfamily
covering the set by the original balls, with multiplicity bounded by
a dimensional constant $N_m$.  Arbitrarily small assigned radii are
allowed.  One proof groups the radii into dyadic bins and, in decreasing
bins, greedily selects uncovered centers.  The selected centers in
each bin are separated by that bin's minimum radius, so only finitely
many are needed in a bounded region.  For the overlap bound at a point,
divide directions into finitely many cones of angular diameter less
than $\pi/3$.  In a cone take the first selected center $x_i$ among
balls containing the point $y$.  For a later such center $x_j$, the
cosine law gives
$|x_i-x_j|<\max(|x_i-y|,|x_j-y|)\le\max(r_i,r_j)$.
Since $x_j$ was not covered by the earlier ball, $|x_i-x_j|>r_i$;
hence $r_j>r_i$.  The decreasing order of the dyadic bins forces
these two radii into the same bin.  Packing in that bin bounds the
number.  A selected center equal to $y$ contributes at most one
additional ball.  This proves the dimensional multiplicity bound.

Choose one good contact and one good ball above each center that has
such a representative, and apply this covering fact.  Since the
original balls lie in $\mathcal E_\lambda$, \eqref{ct:good-ball}
gives
\begin{equation}\label{ct:good-total}
 h^*(E_{\rm good})
 \le A_0p\sum_j\nu(B_j)
 \le A_0N_mp\nu(\mathcal E_\lambda).
\end{equation}
This is the reason for using original balls in this second covering;
the inequality \eqref{ct:good-ball} need not hold on an enlargement.

\paragraph{Choosing the thresholds.}
The measure $\nu$ and the functions $q_v$ are already fixed by the
node.  Choose independent thresholds at the mesh scales with
distribution function
\begin{equation}\label{ct:threshold-law}
 F(t)=
 \begin{cases}
 0,&t\le1,\\
 \exp\bigl(1-1/(t-1)\bigr),&1<t<2,\\
 1,&t\ge2.
 \end{cases}
\end{equation}
For every fixed $K_0>0$, uniformly in real $z$,
\begin{equation}\label{ct:threshold-probability}
 \mathbb P\{|\lambda-z|\le\delta\}
 \le T^{-K_0}+C_{K_0}(\log T)^{-13}F(z-\delta).
\end{equation}
If $F(z+\delta)\le T^{-K_0}$ this is immediate.  Otherwise, put
$a_*=(1+K_0\log T)^{-1}$.  One has $z+\delta-1\ge a_*$,
also when $z+\delta\ge2$, and consequently
$z-\delta-1\ge a_*-2\delta\ge a_*/2$ for large $T$.
Thus this case does not meet the lower endpoint $1$.  On the part
of $[z-\delta,z+\delta]$ below $2$ one has
$(\log F)'=(t-1)^{-2}\le4/a_*^2\le C_{K_0}(\log T)^2$;
above $2$, $\log F=0$.  Integration proves
$F(z+\delta)-F(z-\delta)
 \le C_{K_0}(\log T)^{-13}F(z-\delta)$.
Truncation at $2$ only decreases this difference.  This proves
\eqref{ct:threshold-probability}, including intervals meeting either
endpoint.

Fix $y$, and write $z_k=q_{v_k'}(y)$ and $r_k=F(z_k-\delta)$.
At the $k$th event in \eqref{ct:threshold-event}, all preceding
thresholds have survived, namely $\lambda_{v_j'}\ge z_j-\delta$.
Independence, \eqref{ct:threshold-probability}, and a union bound give
\begin{align*}
 \mathbb P\{y\in\mathcal E_\lambda\}
 &\le KT^{-K_0}
  +C_{K_0}(\log T)^{-13}
       \sum_{k=1}^K r_k\prod_{j<k}(1-r_j)\\
 &\le KT^{-K_0}+C_{K_0}(\log T)^{-13}=o(1).
\end{align*}
The sum telescopes to $1-\prod_k(1-r_k)\le1$; no disjointness of
the candidate events has been assumed.  Since $K=O(T)$, any fixed
$K_0>2$ suffices.  Integrating against the fixed measure $\nu$ and
using Tonelli supplies thresholds with
\[
 \nu(\mathcal E_\lambda)\le o(1)\nu(X).
\]
Choose such thresholds first, and only then construct the contact
measure.  This order makes independence explicit: the contact measure
may depend on the chosen thresholds, but neither $\nu$ nor $q_v$ does.
For large $T$, the last bound pays for $A_0,N_m$ and the fixed factor
between $\nu(X)$ and $\mathcal N_R$.  Equations
\eqref{ct:bad-total} and \eqref{ct:good-total} imply
\[
 h(X)\le\tfrac12p_{i+1}\mathcal N_R+o(\eta)\mathcal N_R.
\]
Because $h(X)=\int w\,\dd L'-\eta\nu(X)$, this is
\eqref{ct:node-conclusion} after increasing the fixed $C$.

Finally solve the recurrence:
\[
 p_0\le2C\eta+2^{-n_0}C_*,\qquad
 2^{-n_0}\le T^{-10N\log2}.
\]
At the initial node $\mathcal N_1\le2d$.  Increasing the fixed $N$
relative to the requested $P$ gives \eqref{ct:separated-estimate}.
There have been at most $n_0$ configuration changes, costing
$(n_0+1)e^b\sqrt d\ll S(T)\sqrt d$ in total.  Thus the fixed root
separation and far cutoff used at every step are valid.  A polynomial
weight or dilation loss is handled by starting with a correspondingly
larger $P$ and using the same Gaussian cell normalization.  This proves
the final assertion as well.
\end{proof}
